\documentclass[aspectratio=169,11pt]{beamer}
\usetheme{Madrid}
\usecolortheme{default}
\setbeamertemplate{navigation symbols}{}
\usepackage[utf8]{inputenc}
\usepackage[T1]{fontenc}
\usepackage{lmodern}
\usepackage[indonesian]{babel}
\usepackage{booktabs}
\usepackage{array}
\usepackage{xcolor}
\usepackage{listings}
\usepackage{amsmath}
\usepackage{amsfonts}
\usepackage{ragged2e}
\usepackage{adjustbox}
\usepackage{multicol}

\definecolor{DeepNavy}{HTML}{0B1220}
\definecolor{Navy}{HTML}{12355B}
\definecolor{Blue}{HTML}{1F77B4}
\definecolor{Teal}{HTML}{00A6A6}
\definecolor{Gold}{HTML}{F2B705}
\definecolor{LightBlue}{HTML}{EAF4FF}
\definecolor{LightTeal}{HTML}{E7FAF8}
\definecolor{SoftGray}{HTML}{F4F6F8}
\definecolor{CodeBg}{HTML}{F8FAFC}
\definecolor{DarkText}{HTML}{15202B}
\definecolor{Warn}{HTML}{D97706}

\setbeamercolor{title}{fg=white}
\setbeamercolor{frametitle}{fg=white,bg=Navy}
\setbeamercolor{structure}{fg=Blue}
\setbeamercolor{normal text}{fg=DarkText,bg=white}
\setbeamercolor{block title}{fg=white,bg=Navy}
\setbeamercolor{block body}{fg=DarkText,bg=LightBlue}
\setbeamercolor{alerted text}{fg=Warn}
\setbeamerfont{title}{size=\LARGE,series=\bfseries}
\setbeamerfont{subtitle}{size=\large}
\setbeamerfont{frametitle}{size=\large,series=\bfseries}
\setbeamerfont{footline}{size=\fontsize{7.2}{8}\selectfont}
\setbeamertemplate{itemize items}[circle]

\lstdefinestyle{py}{
  language=Python,
  basicstyle=\ttfamily\footnotesize,
  keywordstyle=\color{Blue}\bfseries,
  commentstyle=\color{Teal},
  stringstyle=\color{Gold!70!black},
  showstringspaces=false,
  breaklines=true,
  frame=single,
  rulecolor=\color{Navy!35},
  backgroundcolor=\color{CodeBg},
  tabsize=4,
  columns=fullflexible
}
\lstset{style=py}

\newcommand{\mkfooter}{Pemrograman untuk Ilmu Data dan Komputasi}
\newcommand{\prodi}{Program Studi Matematika ITERA}
\setbeamertemplate{footline}{%
\leavevmode%
\hbox{%
\begin{beamercolorbox}[wd=.45\paperwidth,ht=2.35ex,dp=.95ex,leftskip=1.0ex]{author in head/foot}\mkfooter\end{beamercolorbox}%
\begin{beamercolorbox}[wd=.37\paperwidth,ht=2.35ex,dp=.95ex,center]{author in head/foot}\prodi\end{beamercolorbox}%
\begin{beamercolorbox}[wd=.18\paperwidth,ht=2.35ex,dp=.95ex,rightskip=1.0ex plus1fil]{author in head/foot}\hfill\insertframenumber/\inserttotalframenumber\end{beamercolorbox}%
}%
}

\title[\mkfooter]{Pertemuan 7: Aljabar Linier Komputasional}
\subtitle{Vektor, matriks, SPL, residual, eigen, dan data sebagai matriks}
\author{Rifky Fauzi dan Triyana Muliawati, S.Si., M.Si.}
\institute{Program Studi Matematika, Institut Teknologi Sumatera}
\date{Ganjil 2026/2027}

\begin{document}

% 1
\begin{frame}[plain]
\begingroup
\setbeamertemplate{footline}{}
\begin{columns}[T]
\column{.58\textwidth}
\vspace{7mm}
{\color{Navy}\fontsize{22}{25}\selectfont\bfseries Pertemuan 7\\Aljabar Linier Komputasional}\\[4mm]
{\large Pemrograman untuk Ilmu Data dan Komputasi}\\[2mm]
{\normalsize MA25-21008 \quad Kelas RA dan RB}\\[6mm]
{\small Rifky Fauzi\\Triyana Muliawati, S.Si., M.Si.}\\[2mm]
{\small Program Studi Matematika, ITERA\\Semester Ganjil 2026/2027}
\column{.38\textwidth}
\vspace{9mm}
\begin{block}{Fokus hari ini}
\small
\begin{itemize}
  \item membaca vektor dan matriks sebagai struktur data;
  \item menulis algoritma manual sebelum memakai NumPy;
  \item menghitung banyak kerja dari loop;
  \item memakai residual untuk memeriksa hasil numerik.
\end{itemize}
\end{block}
\end{columns}
\endgroup
\end{frame}

% 2
\begin{frame}{Apa yang harus siap untuk LKM M7}
\small
\begin{center}
\begin{tabular}{p{2.2cm}p{5.2cm}p{5.0cm}}
\toprule
Bagian & Yang dikerjakan & Ide algoritma \\
\midrule
M7.1--M7.2 & operasi vektor, norm, element-wise, $A@B$ & loop komponen dan operasi array \\
M7.3--M7.4 & perkalian matriks manual dan transpose & loop bersarang, cek dimensi \\
M7.5--M7.8 & invers 2x2, SPL, residual & determinan, solve, verifikasi $Ax\approx b$ \\
M7.9--M7.10 & eigen dan data sebagai matriks & cek $Av\approx\lambda v$, mean kolom, Gram matrix \\
\bottomrule
\end{tabular}
\end{center}
\vspace{1mm}
\begin{alertblock}{Pola kerja}
Tuliskan rumus matematika, ubah menjadi indeks, baru tulis loop. Setelah itu, jika perlu, bandingkan dengan NumPy.
\end{alertblock}
\end{frame}

% 3
\begin{frame}[fragile]{Vektor: operasi komponen dan norm}
\begin{columns}[T]
\column{.49\textwidth}
\begin{block}{Contoh dari LKM}
\[
x=(1,-2,3),\qquad y=(4,0,-1)
\]
\[
x+y=(5,-2,2),\qquad 2x-y=(-2,-4,7)
\]
\[
x\cdot y=\sum_{i=1}^{n}x_i y_i,
\qquad
\|x\|_2=\sqrt{\sum_{i=1}^{n}x_i^2}
\]
\end{block}
\column{.47\textwidth}
\begin{lstlisting}
def dot(x, y):
    total = 0
    for i in range(len(x)):
        total += x[i] * y[i]
    return total

# norm(x) = sqrt(dot(x,x))
\end{lstlisting}
\begin{alertblock}{Banyak kerja}
Loop berjalan $n$ kali.
\[
T(n)=c_0+c_1 n \quad\Rightarrow\quad O(n)
\]
\end{alertblock}
\end{columns}
\end{frame}

% 4
\begin{frame}[fragile]{Element-wise, matrix product, dan operator \texttt{@}}
\begin{columns}[T]
\column{.45\textwidth}
\[
A=\begin{pmatrix}1&2\\3&4\end{pmatrix},\qquad
B=\begin{pmatrix}2&0\\1&5\end{pmatrix}
\]
\[
A\odot B=\begin{pmatrix}1\cdot2&2\cdot0\\3\cdot1&4\cdot5\end{pmatrix}
=\begin{pmatrix}2&0\\3&20\end{pmatrix}
\]
\[
AB=\begin{pmatrix}
1\cdot2+2\cdot1 & 1\cdot0+2\cdot5\\
3\cdot2+4\cdot1 & 3\cdot0+4\cdot5
\end{pmatrix}
=\begin{pmatrix}4&10\\10&20\end{pmatrix}
\]
\column{.51\textwidth}
\begin{lstlisting}
import numpy as np
A = np.array([[1,2],[3,4]])
B = np.array([[2,0],[1,5]])

A * B   # element-wise
A @ B   # perkalian matriks
\end{lstlisting}
\begin{block}{Perbedaan penting}
\texttt{*} bekerja per elemen dengan posisi yang sama. \texttt{@} memakai baris-kali-kolom.
\end{block}
\end{columns}
\end{frame}

% 5
\begin{frame}[fragile]{Perkalian matriks manual: tiga loop}
Untuk $A\in\mathbb{R}^{m\times p}$ dan $B\in\mathbb{R}^{p\times n}$,
\[
C_{ij}=\sum_{k=1}^{p} A_{ik}B_{kj}.
\]
\begin{columns}[T]
\column{.55\textwidth}
\begin{lstlisting}
def matmul_manual(A, B):
    m, p = len(A), len(A[0])
    p2, n = len(B), len(B[0])
    if p != p2:
        raise ValueError("dimensi tidak cocok")
    C = [[0 for j in range(n)] for i in range(m)]
    for i in range(m):
        for j in range(n):
            for k in range(p):
                C[i][j] += A[i][k] * B[k][j]
    return C
\end{lstlisting}
\column{.41\textwidth}
\begin{alertblock}{Menghitung $T$}
Operasi terdalam berjalan untuk semua pasangan $(i,j,k)$.
\[
T(m,p,n)=m\cdot n\cdot p
\]
Jika matriks persegi $n\times n$:
\[
T(n)=n^3\quad\Rightarrow\quad O(n^3).
\]
\end{alertblock}
\end{columns}
\end{frame}

% 6
\begin{frame}[fragile]{Transpose: menukar peran baris dan kolom}
Jika $A\in\mathbb{R}^{m\times n}$, maka $A^T\in\mathbb{R}^{n\times m}$ dan
\[
(A^T)_{ij}=A_{ji}.
\]
\begin{columns}[T]
\column{.49\textwidth}
\begin{lstlisting}
def transpose(A):
    m, n = len(A), len(A[0])
    T = [[0 for i in range(m)] for j in range(n)]
    for i in range(m):
        for j in range(n):
            T[j][i] = A[i][j]
    return T
\end{lstlisting}
\column{.47\textwidth}
\begin{block}{Contoh}
\[
\begin{pmatrix}1&2&3\\4&5&6\end{pmatrix}^T
=
\begin{pmatrix}1&4\\2&5\\3&6\end{pmatrix}
\]
\end{block}
\begin{alertblock}{Banyak kerja}
Setiap elemen dipindahkan sekali.
\[
T(m,n)=mn\quad\Rightarrow\quad O(mn)
\]
\end{alertblock}
\end{columns}
\end{frame}

% 7
\begin{frame}{Invers 2x2 dan SPL}
\begin{columns}[T]
\column{.50\textwidth}
\begin{block}{Invers matriks 2x2}
\[
A=\begin{pmatrix}a&b\\c&d\end{pmatrix},\qquad
\det(A)=ad-bc
\]
Jika $\det(A)\neq0$,
\[
A^{-1}=\frac{1}{ad-bc}
\begin{pmatrix}d&-b\\-c&a\end{pmatrix}.
\]
Jika $\det(A)=0$, matriks singular dan invers tidak ada.
\end{block}
\column{.46\textwidth}
\begin{block}{SPL sebagai $Ax=b$}
\[
A=\begin{pmatrix}2&1\\5&3\end{pmatrix},\qquad
b=\begin{pmatrix}1\\4\end{pmatrix}
\]
\[
\det(A)=2\cdot3-1\cdot5=1\neq0.
\]
Untuk 2x2, rumus invers manual masih masuk akal. Untuk ukuran besar, gunakan \texttt{np.linalg.solve(A,b)}, bukan membentuk invers terlebih dahulu.
\end{block}
\end{columns}
\end{frame}

% 8
\begin{frame}[fragile]{Residual: cara memeriksa solusi numerik}
Solusi numerik $\hat{x}$ tidak cukup hanya dilihat nilainya. Periksa sisa persamaan:
\[
r=b-A\hat{x},\qquad \|r\|_2.
\]
\begin{columns}[T]
\column{.53\textwidth}
\begin{lstlisting}
def matvec(A, x):
    m, n = len(A), len(A[0])
    y = [0 for i in range(m)]
    for i in range(m):
        for j in range(n):
            y[i] += A[i][j] * x[j]
    return y
\end{lstlisting}
\column{.43\textwidth}
\begin{alertblock}{Banyak kerja}
Matrix-vector product memakai dua loop.
\[
T(m,n)=mn
\]
Menghitung $r=b-A\hat{x}$ dan norm menambah orde $O(m)$.
\[
T(m,n)=mn+m\quad\Rightarrow\quad O(mn)
\]
\end{alertblock}
\begin{block}{Makna residual}
Jika $\|r\|_2$ kecil, maka $\hat{x}$ membuat $A\hat{x}$ dekat dengan $b$.
\end{block}
\end{columns}
\end{frame}

% 9
\begin{frame}{Nilai eigen dan verifikasi}
Untuk matriks kecil, nilai eigen dapat dihitung dari
\[
\det(A-\lambda I)=0.
\]
Contoh:
\[
A=\begin{pmatrix}2&1\\1&2\end{pmatrix},
\quad
A-\lambda I=\begin{pmatrix}2-\lambda&1\\1&2-\lambda\end{pmatrix}
\]
\[
\det(A-\lambda I)=(2-\lambda)^2-1=0
\]
\[
\lambda=1\quad\text{atau}\quad \lambda=3.
\]
\begin{alertblock}{Yang dicek di komputer}
Untuk setiap pasangan $(\lambda,v)$ dari \texttt{np.linalg.eig}, hitung
\[
\|Av-\lambda v\|_2.
\]
Untuk $A\in\mathbb{R}^{n\times n}$, verifikasi $Av$ membutuhkan $O(n^2)$.
\end{alertblock}
\end{frame}

% 10
\begin{frame}{Data sebagai matriks}
Misalkan $X\in\mathbb{R}^{m\times d}$: $m$ observasi dan $d$ fitur.
\[
\bar{x}_j=\frac{1}{m}\sum_{i=1}^{m}X_{ij},
\qquad
(X_c)_{ij}=X_{ij}-\bar{x}_j.
\]
\[
G=X_c^T X_c.
\]
\begin{columns}[T]
\column{.48\textwidth}
\begin{block}{Algoritma}
\begin{enumerate}\small
  \item hitung rata-rata tiap kolom;
  \item kurangi setiap kolom dengan rata-ratanya;
  \item hitung $X_c^T X_c$;
  \item periksa $G$ simetris.
\end{enumerate}
\end{block}
\column{.48\textwidth}
\begin{alertblock}{Banyak kerja}
Mean kolom dan centering:
\[
O(md)
\]
Gram matrix:
\[
G_{jk}=\sum_{i=1}^{m}(X_c)_{ij}(X_c)_{ik}
\]
Ada $d^2$ pasangan fitur, masing-masing menjumlahkan $m$ baris:
\[
O(md^2).
\]
\end{alertblock}
\end{columns}
\end{frame}

% 11
\begin{frame}{Ringkasan algoritma dan orde}
\small
\begin{center}
\begin{tabular}{p{4.4cm}p{4.2cm}p{3.0cm}}
\toprule
Masalah & Operasi utama & Orde \\
\midrule
Dot product, norm vektor & loop semua komponen & $O(n)$ \\
Transpose $m\times n$ & pindahkan tiap elemen & $O(mn)$ \\
Matrix-vector product & dua loop baris-kolom & $O(mn)$ \\
Matrix product $m\times p$ dengan $p\times n$ & tiga loop & $O(mpn)$ \\
Matrix product persegi & tiga loop ukuran $n$ & $O(n^3)$ \\
Verifikasi eigen $Av\approx\lambda v$ & matrix-vector product & $O(n^2)$ \\
Mean dan centering data $m\times d$ & semua entri data & $O(md)$ \\
Gram matrix $X_c^TX_c$ & fitur-pasangan-baris & $O(md^2)$ \\
\bottomrule
\end{tabular}
\end{center}
\begin{alertblock}{Kebiasaan yang dipakai di LKM}
Jawaban numerik perlu disertai cara cek: dimensi cocok, hasil manual kecil, hasil NumPy, dan residual/norm selisih.
\end{alertblock}
\end{frame}

% 12
\begin{frame}{Arah kerja setelah kelas}
\begin{columns}[T]
\column{.50\textwidth}
\begin{block}{Saat mengerjakan LKM M7}
\begin{itemize}
  \item Untuk soal kecil, hitung manual terlebih dahulu.
  \item Untuk list of lists, tulis loop dan cek dimensi.
  \item Untuk NumPy, bedakan \texttt{*}, \texttt{@}, dan \texttt{np.linalg.solve}.
  \item Untuk SPL dan eigen, selalu verifikasi dengan residual atau norm selisih.
\end{itemize}
\end{block}
\column{.46\textwidth}
\begin{block}{Pertanyaan yang harus selalu dijawab}
\begin{itemize}
  \item Ukuran inputnya apa?
  \item Operasi terdalamnya apa?
  \item Berapa kali operasi itu dieksekusi?
  \item Apakah hasilnya bisa diperiksa kembali?
\end{itemize}
\end{block}
\end{columns}
\vspace{1mm}
\begin{center}
\Large rumus $\rightarrow$ indeks $\rightarrow$ loop $\rightarrow$ cek hasil
\end{center}
\end{frame}

\end{document}

